% Mathematical models covered by the Radiant Protocol Master Formula License (RPML) v1.0.

\documentclass[11pt]{article}
\usepackage[margin=1in]{geometry}
\usepackage{amsmath,amssymb,amsthm}
\usepackage{booktabs,hyperref,graphicx}
\usepackage{tikz}
\usetikzlibrary{arrows.meta}

\newtheorem{theorem}{Theorem}[section]
\newtheorem{corollary}[theorem]{Corollary}
\newtheorem{lemma}[theorem]{Lemma}
\newtheorem{remark}[theorem]{Remark}

\title{\textbf{The Meta‑Theorem of Regulatory Dominance over Entropy}\\[4pt]
\small A Single Inequality \(\alpha > 1+\beta\) Governs the Survival of Adaptive Systems}
\author{Anonymous}
\date{}

\begin{document}
\maketitle

\begin{abstract}
Any adaptive system can be characterized by a regulatory capacity $\alpha>0$ and a disturbance amplification factor $\beta\ge0$.
We define the systemic load as
\[
\Lambda(t) \equiv \frac{1+\beta(t)}{\alpha(t)} .
\]
Static stability requires $\Lambda<1$; dynamically, $\Lambda(t)$ remains below the critical threshold iff the relative growth of regulation outpaces that of disturbance:
\[
\boxed{\;\frac{\dot\alpha}{\alpha} > \frac{\dot\beta}{1+\beta}\;} .
\]
From this single inequality we derive the complete phase diagram (stability, criticality, collapse) and, crucially, the formal principle that \textbf{a system dies if and only if its load exceeds its capacity to adapt}.  The result unifies Ashby's requisite variety, the thermodynamics of open systems, Fisher's fundamental theorem, and the stability margins of control engineering.
\end{abstract}

% =================================================================
\section{The Static Condition}
% =================================================================
Consider an adaptive system whose corrective dynamics can be linearized around an equilibrium.
Let $\alpha>0$ be the strength of the negative feedback (the regulatory capacity) and $\beta\ge0$ the factor by which an external disturbance is amplified over one characteristic time.
The system's ability to absorb perturbation is captured by the dimensionless load
\[
\boxed{\;\Lambda(t) = \frac{1+\beta(t)}{\alpha(t)}\;}.
\]
When $\Lambda < 1$, the regulation dominates the amplified disturbance, and the system returns to equilibrium.
When $\Lambda \ge 1$, the disturbance overwhelms the regulation, and any small deviation grows without bound.

\begin{theorem}[Static stability]\label{thm:static}
If $\alpha(t) > 1+\beta(t)$, then the linearised system is globally asymptotically stable.
Conversely, if $\alpha(t) \le 1+\beta(t)$, the system cannot guarantee contraction, and arbitrarily small perturbations may lead to divergence.
\end{theorem}

\section{The Dynamic Extension}
% =================================================================
In realistic environments, both $\alpha$ and $\beta$ evolve.
To avoid crossing the lethal boundary $\Lambda=1$, the load must decrease whenever it is dangerously high.
Differentiating $\ln\Lambda$ gives the condition
\[
\frac{d}{dt}\ln\Lambda(t) = \frac{\dot\beta}{1+\beta} - \frac{\dot\alpha}{\alpha} < 0
\quad\Longleftrightarrow\quad
\boxed{\;\frac{\dot\alpha}{\alpha} > \frac{\dot\beta}{1+\beta}\;} .
\]
This is the fundamental dynamic law: the relative growth of regulation must strictly exceed the relative growth of the amplified disturbance.

\begin{theorem}[Dynamic persistence]\label{thm:dynamic}
Let $\alpha(t)>0$, $\beta(t)\ge0$ be $C^1$ and $\Lambda(0)<1$.
If $\frac{\dot\alpha}{\alpha} > \frac{\dot\beta}{1+\beta}$ holds for all $t\ge0$, then $\Lambda(t)<1$ for all $t\ge0$, and the system remains stable.
\end{theorem}

% =================================================================
\section{The Phase Diagram of Adaptive Systems}
% =================================================================
The relative growth rates define three definitive regimes (Figure~\ref{fig:trajectories}):
\begin{enumerate}
\item \textbf{Stable regime}: $\dot\alpha/\alpha > \dot\beta/(1+\beta)$.  The load decreases; the system is healthy.
\item \textbf{Critical boundary}: $\dot\alpha/\alpha = \dot\beta/(1+\beta)$.  The load is constant; the system is at the edge of collapse.
\item \textbf{Collapse regime}: $\dot\alpha/\alpha < \dot\beta/(1+\beta)$.  The load grows; once $\Lambda\ge 1$, the system dies.
\end{enumerate}

\noindent\textbf{Meta‑Axiom.} From the above we extract the supreme law of adaptive systems:
\[
\boxed{\;\text{A system survives only when its regulatory capacity scales faster than its disturbance.}\;}
\]

\begin{corollary}[Systemic Death Law]\label{cor:death}
A system dies if and only if its systemic load exceeds its regulatory capacity ($\Lambda \ge 1$) and the dynamic condition cannot be restored.
Death is the crossing of $\Lambda = 1$ without a return to the stable regime.
\end{corollary}
This follows directly from the collapse regime: once $\Lambda \ge 1$ and $\dot\Lambda > 0$, the load becomes self‑reinforcing, and the system enters an irreversible entropic cascade.

\begin{corollary}[Zero Legitimacy, Zero Response]\label{cor:zero}
If the regulatory capacity of a system vanishes ($\alpha \to 0$), the systemic load $\Lambda = (1+\beta)/\alpha$ diverges to infinity regardless of the disturbance $\beta$.
Consequently, the system loses any ability to respond to external challenges and dies immediately.
In symbolic terms:
\[
\boxed{\;0\;\text{Legitimacy} \;=\; 0\;\text{Response}\;}.
\]
\end{corollary}

\section{A Concrete Toy Model}
% =================================================================
Let $\alpha(t) = \alpha_0 e^{g_\alpha t}$, $\beta(t) = \beta_0 e^{g_\beta t}$ with $\alpha_0 = 0.8$, $\beta_0 = 0.3$.
Figure~\ref{fig:trajectories} shows two representative trajectories:
\begin{itemize}
\item \textbf{Stable} ($g_\alpha=0.15,\; g_\beta=0.05$): $\Lambda$ remains below $1$.
\item \textbf{Unstable} ($g_\alpha=0.05,\; g_\beta=0.15$): $\Lambda$ crosses $1$ and the system collapses.
\end{itemize}

\begin{figure}[htbp]
\centering
\begin{tikzpicture}[scale=1.05]
  \draw[->] (0,0) -- (12.5,0) node[right] {$t$};
  \draw[->] (0,0) -- (0,5.5) node[above] {$\Lambda(t)$};
  \draw[dashed, gray!60] (0,2.5) -- (12.5,2.5) node[right, black] {Death threshold $1$};
  \draw[blue, thick] plot[domain=0:12, samples=100] (\x, {(1+0.3*exp(0.05*\x))/(0.8*exp(0.15*\x))});
  \draw[red, thick] plot[domain=0:12, samples=100] (\x, {(1+0.3*exp(0.15*\x))/(0.8*exp(0.05*\x))});
  \node[blue] at (9,1.5) {Stable};
  \node[red] at (9,4.3) {Collapse};
\end{tikzpicture}
\caption{Evolution of the systemic load under two growth‑rate regimes.  The red trajectory violates the dynamic inequality and crosses the death threshold.}
\label{fig:trajectories}
\end{figure}

\section{Universality Across Domains}
% =================================================================
The principle described above applies identically to every field where regulation and disturbance interact:

\begin{itemize}
\item \textbf{Control theory}: $\alpha$ is the controller bandwidth, $1+\beta$ the plant amplification.
\item \textbf{Ecology}: $\alpha$ is density‑dependent mortality, $\beta$ the intrinsic growth rate.
\item \textbf{Economics}: $\alpha$ is the stabilisation policy multiplier, $\beta$ the velocity of capital flight.
\item \textbf{Cognitive regulation}: prefrontal inhibitory control ($\alpha$) versus limbic reactivity ($\beta$).
\item \textbf{Institutional viability}: governing capacity ($\alpha$) versus institutional decay/overload ($\beta$).
\end{itemize}
In every case, survival requires $\alpha > 1+\beta$ and the persistent maintenance of $\dot\alpha/\alpha > \dot\beta/(1+\beta)$.

\section{Conclusion}
The Meta‑Theorem of Regulatory Dominance over Entropy provides the universal boundary between the living and the dead for adaptive systems.
The Systemic Load Principle — “a system dies when its load exceeds its capacity to adapt” — is not a metaphor but a precise mathematical consequence of the inequality $\Lambda \ge 1$ accompanied by $\dot\alpha/\alpha \le \dot\beta/(1+\beta)$.
Furthermore, when regulatory capacity completely collapses ($\alpha \to 0$), the system loses all power to respond, formalising the law $0\ \text{Legitimacy} = 0\ \text{Response}$.
This single invariant unifies the core stability conditions of every science of organised complexity.

\begin{thebibliography}{9}
\bibitem{zames1966}
G.~Zames.
\newblock On the input‑output stability of time‑varying nonlinear feedback systems.
\newblock \emph{IEEE Trans. Autom. Control}, 11(2):228‑238, 1966.
\bibitem{samuelson1939}
P.~A.~Samuelson.
\newblock Interactions between the multiplier analysis and the principle of acceleration.
\newblock \emph{Review of Economics and Statistics}, 21(2):75‑78, 1939.
\end{thebibliography}

\end{document}
